\section{Cómo evaluamos cada decisión}\label{sec:evaluacion}

Con 138 parcelas cualquier comparación es ruidosa, así que antes de probar ideas fijamos cómo íbamos a decidir.

\textbf{Error fuera de muestra.} Todas las cifras de RMSE se calculan sobre parcelas que el modelo no vio al
entrenar. Usamos validación cruzada, que reparte las parcelas en 5 grupos (\emph{folds}), entrena con 4 y predice el
quinto, hasta predecir todas. Todo lo que usa el rendimiento, como el rezago espacial, se recalcula dentro de cada
fold para que la parcela evaluada nunca influya en su propia predicción.

\textbf{Dos esquemas.} El esquema A sortea los folds al azar, estratificados por municipio, y repite el sorteo 10
veces. Imita la partición oficial y es el que decide. El esquema B deja fuera zonas completas, uniendo las parcelas a
menos de 2\,km~\cite{roberts2017,ploton2020}, porque la validación al azar tiende a ser optimista cuando hay
autocorrelación espacial. Mide qué tanto depende un modelo de tener vecinos en el entrenamiento.

\textbf{Regla de decisión.} Una idea se acepta si baja el RMSE en al menos 7 de las 10 repeticiones del esquema A y si
pasa una prueba de permutación~\cite{ojala2010}. En esa prueba revolvemos al azar las variables nuevas, dentro de cada
régimen, y repetimos la validación 50 veces con 3 repeticiones cada una. Si el modelo con las variables revueltas
predice igual de bien, la mejora no viene de esas variables y la idea se rechaza. El valor $p$ es la fracción de
permutaciones, contando la real, que predice al menos tan bien como el modelo real. Exigimos $p\le0.05$. Con 50
permutaciones el $p$ más pequeño posible es $1/51=0.0196$, de modo que un valor de 0.02 es el piso de la prueba y no una
medición precisa.

\textbf{Cuántas pruebas hicimos.} Registramos 73 evaluaciones. De ellas, 13 fueron informativas, como las bases de
comparación. Las otras 60 se sometieron a la comparación emparejada por repetición, 10 ganaron en al menos 7
repeticiones y llegaron a la prueba de permutación, y 7 quedaron aceptadas. No corregimos por comparaciones múltiples,
así que los valores $p$ no deben leerse como probabilidades exactas. Con la corrección de Holm solo sobreviven las dos
pruebas de que el régimen supera al estado y al municipio, con $p=0.005$ obtenido con 200 permutaciones. Con el
procedimiento de Benjamini y Hochberg~\cite{benjamini1995}, que controla la tasa de falsos descubrimientos al 5\,\%,
sobreviven cuatro pruebas, las dos anteriores, el rezago espacial y la elección del proceso gaussiano, todas con
$p\le0.0196$. Las tres pruebas con $p=0.039$ no sobreviven. Por esta razón añadimos al final la validación anidada.

\textbf{Controles del resultado final.} Elegir el mejor entre muchos modelos con las mismas particiones infla su
resultado~\cite{cawley2010}. Por eso al final medimos el error esperado con validación anidada~\cite{varma2006},
comparamos modelos con una prueba $t$ corregida por el traslape entre repeticiones~\cite{nadeau2003} y revisamos que
no hubiera fuga de información revolviendo las etiquetas~\cite{kaufman2012}. Estos controles están en la
Sección~\ref{sec:confiabilidad}.
